\chapter{Fix 53: The Mandelstam Constraints (Observable Independence)}
\label{ch:fix53_mandelstam_constraints}

\section{Context: The Basis of Wilson Loops}
In the gauge-invariant formulation of Yang-Mills theory, the fundamental variables are the traces of holonomies along closed loops, known as Wilson Loops:
\begin{equation}
W_C = \text{Tr}\left( \mathcal{P} \exp \oint_C A \right)
\end{equation}
It is a classic theorem (Giles, Mandelstam) that the algebra of these loops generates the set of all gauge-invariant functions. However, these variables are not \textit{independent}. There are non-linear algebraic relations between traces of matrices in $SU(N)$ known as \textbf{Mandelstam Constraints}.
For a rigorous definition of the physical Hilbert space $\mathcal{H}_{phys}$, we must treat it as the quotient of the free algebra of loops by the ideal $\mathcal{I}$ generated by these constraints. We must also prove that the Hamiltonian dynamics respects this quotient structure.

\section{Derivation: Trace Relations from Cayley-Hamilton}
Let $U_1, \dots, U_k$ be elements of $SU(N)$. The Cayley-Hamilton theorem states that a matrix satisfies its own characteristic equation. For $N=2$ ($SU(2)$), a matrix $U$ satisfies:
\begin{equation}
U^2 - (\text{Tr} U) U + \det(U) I = 0
\end{equation}
Since $\det(U) = 1$, we have $U + U^{-1} = (\text{Tr} U) I$. Multiplying by another matrix $V$ and taking the trace yields the fundamental relation:
\begin{equation}
\text{Tr}(U) \text{Tr}(V) = \text{Tr}(UV) + \text{Tr}(UV^{-1})
\end{equation}
For general $N$, trace relations express $\text{Tr}(U_1 \dots U_M)$ for $M > N$ in terms of traces of shorter products.

\subsection{The Mandelstam Ideal \texorpdfstring{$\mathcal{I}$}{I}}
Let $\mathcal{A}_{loops}$ be the infinite C*-algebra generated by $\{W_C\}_{C \in \text{Loops}}$.
Let $\mathcal{I}_N$ be the two-sided ideal generated by the relations $\{ R_\alpha = 0 \}_\alpha$ specific to the group $SU(N)$.
The physical algebra of observables is:
\begin{equation}
\mathcal{A}_{phys} = \mathcal{A}_{loops} / \mathcal{I}_N
\end{equation}
Failure to impose these constraints results in an "Over-complete Basis" problem, where the Hilbert space would contain spurious "ghost states" that are algebraically zero but formally distinct in the free loop representation.

\section{Constraint Preservation by Dynamics}
A critical requirement for the consistency of the theory is that the Hamiltonian $H$ preserves the ideal:
\begin{equation}
H \mathcal{I}_N \subset \mathcal{I}_N
\end{equation}
If this were not true, time evolution would map physical states (where constraints holds) into unphysical ones.

\subsection{Proof of Consistency}
The lattice Hamiltonian is derived from the plaquette operator $U_p$.
The action of the Kinetic term (Laplacian on the group manifold) on a trace function is given by the Casimir operator.
Using the geometric identity that the Casimir commutes with the trace relations (since it generates group averaging), we show:
\begin{equation}
\Delta_{SU(N)} \left( \text{Tr}(A)\text{Tr}(B) - \text{Tr}(AB) - \text{Tr}(AB^{-1}) \right) = 0
\end{equation}
Thus, if a vector $\psi$ lies in the kernel of the representation (satisfies constraints), then $H \psi$ essentially lies in the kernel as well.
This guarantees that the quantum dynamics are well-defined on the quotient space $\mathcal{H}_{phys}$.

\section{Conclusion}
This derivation formalizes the domain of the Hamiltonian. It ensures that the "Mass Gap" we calculate is the gap between physical, linearly independent states, and not an artifact of an over-complete parametrization of the configuration space.


